\chapter{Introducci\'on}
\label{cap:introduccion}

\section{Descripci\'on del sistema}

El sistema administra las inversiones documentadas mediante pagar\'es de inversi\'on con rendimientos mensuales.

La empresa emite los pagar\'es a trav\'es de sus sucursales. Los clientes participan como inversionistas (acreedores) y, cuando existe, un aval garantiza el pago de las obligaciones de la empresa. El personal de la empresa es el \'unico que opera el sistema; los clientes y los avales no acceden directamente a la base de datos, sino que reciben de la empresa la informaci\'on que les corresponde. El pagar\'e establece los siguientes elementos:

\begin{itemize}[leftmargin=*]
	\item \textbf{Datos del pagar\'e:} n\'umero de pagar\'e, monto de la inversi\'on y su moneda, y el lugar y la fecha de suscripci\'on.
	\item \textbf{Partes:} el emisor y el inversionista, cada uno con su nombre completo y domicilio.
	\item \textbf{Fecha de vencimiento:} fecha en la que el emisor debe pagar la cantidad principal de la inversi\'on.
	\item \textbf{Intereses ordinarios y pagos mensuales:} el emisor paga un inter\'es ordinario mensual a una tasa fija (o una tasa anual divisible entre doce), calculado sobre el saldo del capital insoluto. Los pagos se realizan de manera mensual y consecutiva en un d\'ia espec\'ifico de cada mes, y concluyen en la fecha de vencimiento junto con la devoluci\'on del capital principal.
	\item \textbf{Intereses moratorios:} en caso de mora o retraso en el pago de cualquiera de las mensualidades de intereses o del capital principal, el emisor paga un inter\'es moratorio a una tasa mensual sobre los saldos vencidos y no pagados, desde la fecha de su vencimiento hasta su total liquidaci\'on.
	\item \textbf{Disposiciones generales:} el pagar\'e se rige por la Ley General de T\'itulos y Operaciones de Cr\'edito, y las partes se someten a la jurisdicci\'on de los tribunales de una ciudad determinada.
	\item \textbf{Aval (opcional):} persona que garantiza el pago de las obligaciones del emisor, con su nombre completo, domicilio e identificaci\'on.
\end{itemize}

\section{Objetivo general}

Generar un esquema que permita tener un control adecuado de un sistema de inversiones, as\'i como la optimizaci\'on de las liquidaciones de las inversiones proporcionadas por los clientes. La base de datos deber\'a almacenar de manera adecuada los datos de los clientes, asegurando su integridad.

\section{Objetivos espec\'ificos}

\begin{enumerate}[leftmargin=*]
	\item Consultar los datos de las personas que tengan inversiones otorgadas en un periodo, por nombre.
	\item Consultar los datos de las inversiones que est\'en vencidas en un periodo.
	\item Consultar los estados de cuenta de los inversionistas.
	\item Consultar la inversi\'on y su tabla de inter\'es en un periodo.
	\item Consultar las inversiones hechas en un periodo.
\end{enumerate}

\section{Consultas de los objetivos espec\'ificos}
\label{sec:consultas-objetivos}

Cada objetivo espec\'ifico se resuelve con una consulta que puede ejecutar cualquier cuenta de operaci\'on (gerente, ejecutivo, tesorer\'ia o auditor), ya que solo requiere el privilegio \nem{SELECT}. Cada objetivo se acompa\~na de una captura que muestra el formato de la respuesta, igual en ambos gestores, obtenida con la cuenta \nem{op\us gerente}. Las consultas solo leen datos, por lo que no requieren identificarse con \nem{SET} (secci\'on~\ref{sec:identificacion}).

\paragraph{Par\'ametros.} Los valores de entrada (el periodo, el nombre buscado, el cliente o el pagar\'e) se escriben directamente en el \nem{WHERE} de cada consulta. Para hacer otra consulta solo se cambian esos valores.

\paragraph{Versiones por instancia.} Los objetivos 1 y 4 se resuelven con la misma consulta en las dos instancias. Los objetivos 2, 3 y 5 necesitan datos calculados, como el estado de una inversi\'on, el inter\'es pendiente o la mora estimada, y tienen una versi\'on por instancia:

\begin{itemize}[leftmargin=*]
	\item \textbf{Primera instancia:} la consulta se resuelve directamente sobre las tablas, por lo que incluye los c\'alculos. Es m\'as larga: obtiene primero el saldo o lo pagado de cada periodo de inter\'es en una expresi\'on com\'un (\nem{periodo}) y despu\'es lo agrega por inversi\'on.
	\item \textbf{Segunda instancia:} la consulta se resuelve sobre las vistas derivadas de la segmentaci\'on, \nem{interes\us saldo} e \nem{inversion\us estado} (secci\'on~\ref{sec:segmentacion-derivada}), que ya contienen esos c\'alculos. La consulta se reduce a elegir columnas y filtrar.
\end{itemize}

Con los datos de prueba, ambas versiones devuelven exactamente el mismo resultado.

\paragraph{Fecha de corte.} El estado de una inversi\'on y los montos vencidos dependen del d\'ia en que se consultan. Las consultas calculan siempre a la fecha del sistema (\nem{CURRENT\us DATE}), por lo que el mismo dato produce resultados distintos conforme pasa el tiempo. Las capturas y los ejemplos de este manual se obtuvieron con los datos de prueba, cuyo \'ultimo movimiento es del 28 de octubre de 2026; los valores que dependen de la fecha, como los d\'ias vencida y la mora estimada, cambian cada d\'ia.

\paragraph{Criterios comunes.} Las consultas aplican las reglas de negocio del modelo conceptual (cap\'itulo~\ref{cap:conceptual}):

\begin{itemize}[leftmargin=*]
	\item El estado de una inversi\'on es \textit{liquidada} si tiene una liquidaci\'on registrada, \textit{vencida} si su fecha de vencimiento ya pas\'o y \textit{vigente} en los dem\'as casos.
	\item El inter\'es de un periodo es exigible cuando su fecha final es anterior o igual a hoy. Se consideran todos los pagos registrados.
	\item El inter\'es moratorio se calcula como en los datos de prueba: sobre la tasa moratoria mensual, en proporci\'on a los d\'ias de atraso y con meses de 30 d\'ias.
\end{itemize}

\paragraph{Diferencias entre gestores.} Solo cambian de un gestor a otro la b\'usqueda por nombre del objetivo~1, la presentaci\'on en formato vertical de los datos del pagar\'e del objetivo~4 y la resta de fechas, que en la primera instancia solo aparece en el objetivo~2 (en PostgreSQL restar dos fechas devuelve los d\'ias; en MySQL se utiliza \nem{DATEDIFF}). Las dem\'as consultas se ejecutan sin cambios en ambos gestores.

\subsection{Objetivo 1}

\textit{Consultar los datos de las personas que tengan inversiones otorgadas en un periodo, por nombre.}

La consulta lista a los inversionistas que suscribieron al menos un pagar\'e en el periodo y cuyo nombre completo contiene el texto buscado. Devuelve un rengl\'on por persona, con sus datos, el n\'umero de inversiones del periodo y el capital total invertido. Es la misma en las dos instancias.

\paragraph{B\'usqueda sin acentos.} La intercalaci\'on de la base de datos en MySQL compara los textos sin distinguir may\'usculas, min\'usculas ni acentos (cap\'itulo~\ref{cap:implementacion}). PostgreSQL s\'i los distingue, por lo que su consulta usa \nem{ILIKE} para ignorar las may\'usculas y la funci\'on \nem{translate} para sustituir las vocales acentuadas y la letra \~n del nombre por su letra base.

\paragraph{PostgreSQL.}

\begin{verbatim}
	SELECT c.cliente_num, c.cliente_rfc,
	       c.cliente_nombre, c.cliente_apellido,
	       c.cliente_direccion, c.cliente_telefono,
	       COUNT(*)                 AS inversiones,
	       SUM(i.inversion_capital) AS capital_total
	FROM inversion i
	  JOIN cliente c ON c.cliente_num = i.cliente_num
	WHERE i.pagare_fecha BETWEEN '2025-01-01' AND '2025-12-31'
	  AND translate(c.cliente_nombre || ' ' || c.cliente_apellido,
	                'ÁÉÍÓÚÜÑáéíóúüñ', 'AEIOUUNaeiouun') ILIKE '%jose%'
	GROUP BY c.cliente_num
	ORDER BY c.cliente_apellido, c.cliente_nombre;
\end{verbatim}

\paragraph{MySQL.}

\begin{verbatim}
	SELECT c.cliente_num, c.cliente_rfc,
	       c.cliente_nombre, c.cliente_apellido,
	       c.cliente_direccion, c.cliente_telefono,
	       COUNT(*)                 AS inversiones,
	       SUM(i.inversion_capital) AS capital_total
	FROM inversion i
	  JOIN cliente c ON c.cliente_num = i.cliente_num
	WHERE i.pagare_fecha BETWEEN '2025-01-01' AND '2025-12-31'
	  AND CONCAT(c.cliente_nombre, ' ', c.cliente_apellido) LIKE '%jose%'
	GROUP BY c.cliente_num
	ORDER BY c.cliente_apellido, c.cliente_nombre;
\end{verbatim}

Con los par\'ametros del ejemplo, el a\~no 2025 y la b\'usqueda \nem{jose}, ambas consultas devuelven un rengl\'on: el cliente 1, con una inversi\'on de \$75{,}000.00.

\captura{objetivo1.png}{Formato de la respuesta del objetivo 1}{objetivo1}

\subsection{Objetivo 2}

\textit{Consultar los datos de las inversiones que est\'en vencidas en un periodo.}

La consulta lista las inversiones vencidas, es decir, las que ya pasaron su fecha de vencimiento sin tener una liquidaci\'on registrada, cuyo vencimiento cae dentro del periodo. Para cada una muestra los d\'ias vencida, los periodos sin pagar, el inter\'es vencido y la mora estimada:

\begin{itemize}[leftmargin=*]
	\item \textbf{D\'ias vencida:} d\'ias transcurridos desde el vencimiento hasta hoy.
	\item \textbf{Periodos sin pagar e inter\'es vencido:} periodos terminados que conservan un saldo, y la suma de esos saldos. Un periodo con pago parcial cuenta por su saldo restante.
	\item \textbf{Mora estimada:} inter\'es moratorio que la empresa adeudar\'ia si liquidara hoy. Corresponde al capital por los d\'ias desde el vencimiento, m\'as el saldo de cada periodo por los d\'ias desde su fecha final. Es un c\'alculo y no un valor almacenado: se registra hasta que la empresa lo paga, en \nem{pago} o en \nem{liquidacion}.
\end{itemize}

\paragraph{Primera instancia, PostgreSQL.}

\begin{verbatim}
	WITH periodo AS (
	  -- Saldo de interés de cada periodo ya terminado
	  SELECT t.inversion_num, t.interes_fecha_final,
	         t.interes_monto - COALESCE(SUM(g.pago_interes), 0) AS saldo
	  FROM interes t
	    LEFT JOIN pago g
	      ON g.inversion_num = t.inversion_num
	     AND g.interes_num = t.interes_num
	  WHERE t.interes_fecha_final < CURRENT_DATE
	  GROUP BY t.inversion_num, t.interes_num,
	           t.interes_fecha_final, t.interes_monto
	)
	SELECT i.inversion_num, i.pagare_num,
	       CONCAT(c.cliente_nombre, ' ', c.cliente_apellido) AS cliente,
	       s.sucursal_nombre,
	       i.inversion_capital AS capital,
	       i.inversion_fecha_vencimiento AS vencimiento,
	       CURRENT_DATE - i.inversion_fecha_vencimiento AS dias_vencida,
	       SUM(CASE WHEN p.saldo > 0 THEN 1 ELSE 0 END) AS periodos_sin_pagar,
	       SUM(CASE WHEN p.saldo > 0 THEN p.saldo ELSE 0 END) AS interes_vencido,
	       ROUND(i.inversion_capital * i.inversion_tasa_moratoria / 100
	             * (CURRENT_DATE - i.inversion_fecha_vencimiento) / 30
	             + SUM(CASE WHEN p.saldo > 0
	                        THEN p.saldo * i.inversion_tasa_moratoria / 100
	                             * (CURRENT_DATE - p.interes_fecha_final) / 30
	                        ELSE 0 END), 2) AS mora_estimada
	FROM inversion i
	  JOIN cliente c  ON c.cliente_num = i.cliente_num
	  JOIN sucursal s ON s.sucursal_num = i.sucursal_num
	  JOIN periodo p  ON p.inversion_num = i.inversion_num
	WHERE i.inversion_fecha_vencimiento BETWEEN '2025-01-01' AND '2025-12-31'
	  AND i.inversion_fecha_vencimiento < CURRENT_DATE
	  AND NOT EXISTS (SELECT 1 FROM liquidacion l
	                  WHERE l.inversion_num = i.inversion_num)
	GROUP BY i.inversion_num, c.cliente_num, s.sucursal_num
	ORDER BY i.inversion_fecha_vencimiento;
\end{verbatim}

\paragraph{Primera instancia, MySQL.} Es igual a la de PostgreSQL, salvo que cada resta de fechas usa \nem{DATEDIFF}.

\begin{verbatim}
	WITH periodo AS (
	  -- Saldo de interés de cada periodo ya terminado
	  SELECT t.inversion_num, t.interes_fecha_final,
	         t.interes_monto - COALESCE(SUM(g.pago_interes), 0) AS saldo
	  FROM interes t
	    LEFT JOIN pago g
	      ON g.inversion_num = t.inversion_num
	     AND g.interes_num = t.interes_num
	  WHERE t.interes_fecha_final < CURRENT_DATE
	  GROUP BY t.inversion_num, t.interes_num,
	           t.interes_fecha_final, t.interes_monto
	)
	SELECT i.inversion_num, i.pagare_num,
	       CONCAT(c.cliente_nombre, ' ', c.cliente_apellido) AS cliente,
	       s.sucursal_nombre,
	       i.inversion_capital AS capital,
	       i.inversion_fecha_vencimiento AS vencimiento,
	       DATEDIFF(CURRENT_DATE, i.inversion_fecha_vencimiento) AS dias_vencida,
	       SUM(CASE WHEN p.saldo > 0 THEN 1 ELSE 0 END) AS periodos_sin_pagar,
	       SUM(CASE WHEN p.saldo > 0 THEN p.saldo ELSE 0 END) AS interes_vencido,
	       ROUND(i.inversion_capital * i.inversion_tasa_moratoria / 100
	             * DATEDIFF(CURRENT_DATE, i.inversion_fecha_vencimiento) / 30
	             + SUM(CASE WHEN p.saldo > 0
	                        THEN p.saldo * i.inversion_tasa_moratoria / 100
	                             * DATEDIFF(CURRENT_DATE, p.interes_fecha_final) / 30
	                        ELSE 0 END), 2) AS mora_estimada
	FROM inversion i
	  JOIN cliente c  ON c.cliente_num = i.cliente_num
	  JOIN sucursal s ON s.sucursal_num = i.sucursal_num
	  JOIN periodo p  ON p.inversion_num = i.inversion_num
	WHERE i.inversion_fecha_vencimiento BETWEEN '2025-01-01' AND '2025-12-31'
	  AND i.inversion_fecha_vencimiento < CURRENT_DATE
	  AND NOT EXISTS (SELECT 1 FROM liquidacion l
	                  WHERE l.inversion_num = i.inversion_num)
	GROUP BY i.inversion_num, c.cliente_num, s.sucursal_num
	ORDER BY i.inversion_fecha_vencimiento;
\end{verbatim}

\paragraph{Segunda instancia.} Igual en ambos gestores. Las columnas provienen de la vista \nem{inversion\us estado} (tabla~\ref{tab:vista-estado}).

\begin{verbatim}
	SELECT inversion_num, pagare_num, cliente, sucursal_nombre, capital,
	       vencimiento, dias_vencida, periodos_sin_pagar,
	       interes_vencido, mora_estimada
	FROM inversion_estado
	WHERE estado = 'Vencida'
	  AND vencimiento BETWEEN '2025-01-01' AND '2025-12-31'
	ORDER BY vencimiento;
\end{verbatim}

Con el periodo del ejemplo, el a\~no 2025, ambas versiones devuelven las inversiones 16, 18, 20 y 17, en ese orden de vencimiento, cada una con dos periodos sin pagar. La mora estimada de la inversi\'on 16 es el capital de \$500{,}000 por 2.5\,\% mensual y los d\'ias transcurridos desde el 5 de julio de 2025 entre 30, m\'as el saldo de \$4{,}250 de cada periodo pendiente por la misma tasa y sus d\'ias de atraso; crece cada d\'ia.

\captura{objetivo2.png}{Formato de la respuesta del objetivo 2}{objetivo2}

\subsection{Objetivo 3}

\textit{Consultar los estados de cuenta de los inversionistas.}

El estado de cuenta de un inversionista, identificado por su RFC, muestra un rengl\'on por cada una de sus inversiones, con:

\begin{itemize}[leftmargin=*]
	\item El estado de la inversi\'on.
	\item El inter\'es exigible, el pagado y el pendiente.
	\item El inter\'es moratorio pagado, tanto en los pagos mensuales como en la liquidaci\'on.
	\item El capital devuelto.
\end{itemize}

Si se unieran directamente \nem{interes} y \nem{pago}, el monto de cada periodo se repetir\'ia por cada pago y las sumas resultar\'ian multiplicadas. Por ello, la versi\'on de la primera instancia obtiene antes lo pagado de cada periodo (\nem{periodo}). Ambas versiones son iguales en los dos gestores.

\paragraph{Primera instancia.}

\begin{verbatim}
	WITH periodo AS (
	  -- Lo pagado de cada periodo de interés
	  SELECT t.inversion_num, t.interes_fecha_final, t.interes_monto,
	         COALESCE(SUM(g.pago_interes), 0) AS pagado,
	         COALESCE(SUM(g.pago_mora), 0)    AS mora_pagada
	  FROM interes t
	    LEFT JOIN pago g
	      ON g.inversion_num = t.inversion_num
	     AND g.interes_num = t.interes_num
	  GROUP BY t.inversion_num, t.interes_num,
	           t.interes_fecha_final, t.interes_monto
	)
	SELECT i.inversion_num, i.pagare_num, i.pagare_fecha AS suscripcion,
	       i.inversion_fecha_vencimiento AS vencimiento,
	       i.inversion_capital AS capital,
	       CASE WHEN l.inversion_num IS NOT NULL THEN 'Liquidada'
	            WHEN i.inversion_fecha_vencimiento < CURRENT_DATE THEN 'Vencida'
	            ELSE 'Vigente'
	       END AS estado,
	       SUM(CASE WHEN p.interes_fecha_final <= CURRENT_DATE
	                THEN p.interes_monto ELSE 0 END) AS interes_exigible,
	       SUM(p.pagado) AS interes_pagado,
	       SUM(CASE WHEN p.interes_fecha_final <= CURRENT_DATE
	                THEN p.interes_monto ELSE 0 END) - SUM(p.pagado)
	         AS interes_pendiente,
	       SUM(p.mora_pagada) + COALESCE(l.liquidacion_mora, 0) AS mora_pagada,
	       COALESCE(l.liquidacion_capital, 0.00) AS capital_devuelto
	FROM cliente c
	  JOIN inversion i ON i.cliente_num = c.cliente_num
	  JOIN periodo p   ON p.inversion_num = i.inversion_num
	  LEFT JOIN liquidacion l ON l.inversion_num = i.inversion_num
	WHERE c.cliente_rfc = 'MEAP850304BV0'
	GROUP BY i.inversion_num, l.inversion_num
	ORDER BY i.pagare_fecha;
\end{verbatim}

\paragraph{Segunda instancia.} Todos los datos provienen de la vista \nem{inversion\us estado}.

\begin{verbatim}
	SELECT inversion_num, pagare_num, suscripcion, vencimiento, capital,
	       estado, interes_exigible, interes_pagado, interes_pendiente,
	       mora_pagada, capital_devuelto
	FROM inversion_estado
	WHERE cliente_rfc = 'MEAP850304BV0'
	ORDER BY suscripcion;
\end{verbatim}

El ejemplo corresponde a la clienta Patricia Mendoza Aguilar, que tiene dos inversiones. La inversi\'on 10 est\'a liquidada: se devolvi\'o con 20 d\'ias de atraso, con un inter\'es moratorio de $75{,}000 \times 2.5\mkern 3mu\% \times 20/30 = 1{,}250.00$. La inversi\'on 35 est\'a vigente y al corriente; su mora pagada de \$7.50 proviene del pago parcial del periodo de junio de 2026.

\captura{objetivo3.png}{Formato de la respuesta del objetivo 3}{objetivo3}

\subsection{Objetivo 4}

\textit{Consultar la inversi\'on y su tabla de inter\'es en un periodo.}

El objetivo se resuelve con dos consultas, que reciben el n\'umero de pagar\'e, y son las mismas en las dos instancias. La primera muestra los datos de la inversi\'on, que corresponden a los elementos del pagar\'e. Como devuelve un solo rengl\'on con muchas columnas, se presenta en formato vertical: en PostgreSQL con el comando \nem{\textbackslash x} de \nem{psql}, y en MySQL terminando la sentencia con \nem{\textbackslash G} en lugar del punto y coma. Cuando la inversi\'on no tiene aval, los datos del aval aparecen como \nem{NULL}.

\paragraph{PostgreSQL.}

\begin{verbatim}
	\x on
	SELECT i.inversion_num, i.pagare_num, i.pagare_fecha,
	       s.sucursal_ciudad AS lugar_suscripcion,
	       c.cliente_nombre, c.cliente_apellido,
	       a.aval_nombre, a.aval_apellido,
	       i.inversion_capital, i.inversion_moneda, i.inversion_tasa,
	       i.inversion_tasa_moratoria, i.inversion_dia_pago,
	       i.inversion_fecha_vencimiento, i.pagare_jurisdiccion
	FROM inversion i
	  JOIN sucursal s ON s.sucursal_num = i.sucursal_num
	  JOIN cliente c  ON c.cliente_num = i.cliente_num
	  LEFT JOIN aval a ON a.aval_num = i.aval_num
	WHERE i.pagare_num = '0018/2024';
	\x off
\end{verbatim}

\paragraph{MySQL.}

\begin{verbatim}
	SELECT i.inversion_num, i.pagare_num, i.pagare_fecha,
	       s.sucursal_ciudad AS lugar_suscripcion,
	       c.cliente_nombre, c.cliente_apellido,
	       a.aval_nombre, a.aval_apellido,
	       i.inversion_capital, i.inversion_moneda, i.inversion_tasa,
	       i.inversion_tasa_moratoria, i.inversion_dia_pago,
	       i.inversion_fecha_vencimiento, i.pagare_jurisdiccion
	FROM inversion i
	  JOIN sucursal s ON s.sucursal_num = i.sucursal_num
	  JOIN cliente c  ON c.cliente_num = i.cliente_num
	  LEFT JOIN aval a ON a.aval_num = i.aval_num
	WHERE i.pagare_num = '0018/2024'\G
\end{verbatim}

La segunda consulta, igual en ambos gestores, muestra los periodos de inter\'es cuya fecha final cae dentro del periodo consultado. Cada periodo se identifica por su n\'umero (\nem{interes\us num}) dentro de la inversi\'on, que es tambi\'en la columna con la que se une a sus pagos. Para cada uno indica el n\'umero de pagos, la fecha del \'ultimo, el monto pagado y el saldo. Un periodo est\'a pagado cuando su saldo es cero; si el \'ultimo pago es posterior a su fecha final, se pag\'o con atraso, y si no tiene pagos y su fecha final ya pas\'o, est\'a vencido.

\begin{verbatim}
	SELECT t.interes_num AS periodo,
	       t.interes_fecha_inicial AS inicio, t.interes_fecha_final AS fin,
	       t.interes_monto AS interes,
	       COUNT(g.pago_fecha) AS pagos,
	       MAX(g.pago_fecha)   AS ultimo_pago,
	       COALESCE(SUM(g.pago_interes), 0) AS pagado,
	       t.interes_monto - COALESCE(SUM(g.pago_interes), 0) AS saldo
	FROM inversion i
	  JOIN interes t ON t.inversion_num = i.inversion_num
	  LEFT JOIN pago g
	    ON g.inversion_num = t.inversion_num
	   AND g.interes_num = t.interes_num
	WHERE i.pagare_num = '0018/2024'
	  AND t.interes_fecha_final BETWEEN '2025-06-01' AND '2026-03-31'
	GROUP BY t.inversion_num, t.interes_num, t.interes_fecha_inicial,
	         t.interes_fecha_final, t.interes_monto
	ORDER BY t.interes_num;
\end{verbatim}

El ejemplo corresponde al pagar\'e \nem{0018/2024} (inversi\'on 21), que no tiene aval y est\'a vencido. Entre junio de 2025 y marzo de 2026 tiene diez periodos. El primero, que vence el 22 de junio de 2025, se pag\'o en dos partes y el \'ultimo pago es posterior a su fecha final, por lo que se pag\'o con atraso. Los siete siguientes est\'an pagados, y los dos \'ultimos, vencidos sin pago.

\captura{objetivo4.png}{Formato de las respuestas del objetivo 4}{objetivo4}

\subsection{Objetivo 5}

\textit{Consultar las inversiones hechas en un periodo.}

La consulta lista los pagar\'es suscritos en el periodo, con la sucursal, el cliente, el capital, la tasa mensual, el plazo, la fecha de vencimiento y el estado de cada inversi\'on. Es igual en ambos gestores.

\paragraph{Primera instancia.} El plazo se obtiene contando los periodos de su tabla de inter\'es, ya que cada inversi\'on tiene un periodo por cada mes de plazo, y el estado se determina con un \nem{CASE}.

\begin{verbatim}
	SELECT i.inversion_num, i.pagare_num, i.pagare_fecha AS suscripcion,
	       s.sucursal_nombre,
	       CONCAT(c.cliente_nombre, ' ', c.cliente_apellido) AS cliente,
	       i.inversion_capital AS capital, i.inversion_tasa AS tasa,
	       (SELECT COUNT(*) FROM interes t
	        WHERE t.inversion_num = i.inversion_num) AS plazo_meses,
	       i.inversion_fecha_vencimiento AS vencimiento,
	       CASE WHEN l.inversion_num IS NOT NULL THEN 'Liquidada'
	            WHEN i.inversion_fecha_vencimiento < CURRENT_DATE THEN 'Vencida'
	            ELSE 'Vigente'
	       END AS estado
	FROM inversion i
	  JOIN sucursal s ON s.sucursal_num = i.sucursal_num
	  JOIN cliente c  ON c.cliente_num = i.cliente_num
	  LEFT JOIN liquidacion l ON l.inversion_num = i.inversion_num
	WHERE i.pagare_fecha BETWEEN '2026-01-01' AND '2026-10-28'
	ORDER BY i.pagare_fecha;
\end{verbatim}

\paragraph{Segunda instancia.} El plazo y el estado provienen de la vista \nem{inversion\us estado}.

\begin{verbatim}
	SELECT inversion_num, pagare_num, suscripcion, sucursal_nombre, cliente,
	       capital, tasa, plazo_meses, vencimiento, estado
	FROM inversion_estado
	WHERE suscripcion BETWEEN '2026-01-01' AND '2026-10-28'
	ORDER BY suscripcion;
\end{verbatim}

Del 1 de enero al 28 de octubre de 2026 se suscribieron cinco pagar\'es, las inversiones 36 a 40.

\captura{objetivo5.png}{Formato de la respuesta del objetivo 5}{objetivo5}
